%============================ IMPLEMENTACIÓN (ENTREGA 2) ====================
\chapter{Implementación del prototipo}
\label{cap:implementacion}

Este capítulo documenta la implementación real del sistema. El proyecto suma más de 3\,900 líneas de código distribuidas entre el firmware del ESP32 (777 líneas), el servidor Node.js (665 líneas más 274 en los módulos \code{lib/}) y el frontend (913 líneas de lógica, 419 de estructura y 1\,272 de estilos). Se presentan los fragmentos más representativos; los archivos completos se referencian en los anexos y se acompañan de capturas donde procede.

\section{Programación del firmware (Arduino / ESP32)}
El firmware se desarrolló en C++ sobre el \emph{core} de Arduino para ESP32, empleando las librerías \code{WiFi}, \code{WebServer}, \code{DNSServer}, \code{WebSocketsClient} (Markus Sattler), la pila BLE nativa y \code{Adafruit\_SSD1306} para la pantalla. El bucle principal coordina el escaneo, el \emph{heartbeat}, la lectura del pulsador y, en modo SoftAP, la atención del portal cautivo.

\subsection{Estructura de modos y bucle principal}
El Listado~\ref{lst:loop} muestra el despacho por modo dentro de \code{loop()}.

\begin{lstlisting}[style=cpp,caption={Bucle principal y despacho por modo (firmware)},label={lst:loop}]
void loop() {
  webSocket.loop();
  button.loop();

  // Heartbeat cada 3 s hacia el servidor
  if (isWsConnected && (millis() - lastHeartbeatTime > 3000)) {
    webSocket.sendTXT("{\"action\":\"heartbeat\"" + tokenField() + ...);
    lastHeartbeatTime = millis();
  }
  if (button.isPressed()) applyMode((current_mode + 1) % 4); // pulsador D15

  switch (current_mode) {
    case 0: if (millis() - lastScanTime > 3000) { scanAndAnalyzeSpectrum(); lastScanTime = millis(); } break;
    case 1: if (millis() - lastScanTime > 3000) { scanAndStreamBLE();      lastScanTime = millis(); } break;
    case 2: if (captiveRunning) { dnsServer.processNextRequest(); captiveServer.handleClient(); } break;
    case 3: break; // standby
  }
}
\end{lstlisting}

El cambio de modo se realiza físicamente con el pulsador (GPIO15); en cada
pulsación el firmware avanza al siguiente modo y refresca la pantalla OLED con la
información propia de ese estado. Las Figuras siguientes muestran el dispositivo
físico operando en los modos de escaneo Wi\nobreakdash-Fi y BLE.

\captura[0.7]{foto-oled-wifi}{Pantalla OLED en modo escaneo Wi\nobreakdash-Fi (modo 0): total de redes, canal más saturado y canal recomendado.}

\captura[0.7]{foto-oled-ble}{Pantalla OLED en modo escaneo BLE (modo 1): dispositivos detectados y RSSI del más cercano.}

\subsection{Escaneo de espectro Wi-Fi}
En el modo 0, el firmware realiza un escaneo asíncrono, cuenta los puntos de acceso por canal y construye una trama JSON con los dispositivos y el histograma de canales, que envía por WebSocket (Listado~\ref{lst:scan}).

\begin{lstlisting}[style=cpp,caption={Escaneo y serialización del espectro Wi-Fi},label={lst:scan}]
String json = "{\"type\":\"spectrum_wifi\"" + tokenField() +
              ",\"deviceId\":\"ESP32-EILOR\",\"devices\":[";
for (int i = 0; i < n; ++i) {
  int ch = WiFi.channel(i);
  if (ch >= 1 && ch <= 13) channelCounts[ch]++;
  if (i > 0) json += ",";
  json += "{\"ssid\":\""   + WiFi.SSID(i)     + "\","
          "\"bssid\":\""   + WiFi.BSSIDstr(i) + "\","
          "\"rssi\":"      + String(WiFi.RSSI(i))    + ","
          "\"channel\":"   + String(ch)              + "}";
}
json += "],\"channels\":{ ... }}";
if (isWsConnected) webSocket.sendTXT(json);
\end{lstlisting}

\subsection{Portal cautivo y \emph{buffer} offline}
Al entrar en modo SoftAP, \code{startCaptivePortal()} inicia un servidor DNS que redirige todas las consultas a la IP del ESP32 y un servidor HTTP que sirve el formulario. El manejador del formulario delega en \code{queueOrSendReg()}, que envía el registro o lo encola si no hay enlace (Listado~\ref{lst:buffer}).

\begin{lstlisting}[style=cpp,caption={Envío con buffer offline (firmware)},label={lst:buffer}]
void queueOrSendReg(const String& nombre, const String& apellido,
                    const String& telefono, const String& apSsid) {
  if (isWsConnected) {
    String j = buildRegJson(nombre, apellido, telefono, apSsid); // local (sendTXT pide String&)
    webSocket.sendTXT(j);
  } else if (pendingCount < MAX_PENDING) {          // sin enlace: encolar
    pendingRegs[pendingCount++] = { nombre, apellido, telefono, apSsid };
  }
}
void flushPendingRegs() {                            // al reconectar: reenviar
  if (!isWsConnected || pendingCount == 0) return;
  for (int i = 0; i < pendingCount; i++) {
    String j = buildRegJson(pendingRegs[i].nombre, pendingRegs[i].apellido,
                            pendingRegs[i].telefono, pendingRegs[i].apSsid);
    webSocket.sendTXT(j); delay(30);
  }
  pendingCount = 0;
}
\end{lstlisting}

\captura[0.7]{foto-oled-softap}{Pantalla OLED en modo SoftAP: SSID emitido, canal, IP del portal (192.168.4.1) y contadores de clientes y registros.}

\section{Implementación del servidor (Node.js)}
\subsection{Módulo de análisis de canales}
El corazón analítico es \code{analyzeChannels()}, que calcula un puntaje de interferencia por canal ponderando el solapamiento de canales adyacentes (hasta 4 canales de distancia) y la intensidad de señal, y recomienda el canal más limpio del conjunto no solapado (Listado~\ref{lst:analysis}).

\begin{lstlisting}[style=js,caption={Puntaje de interferencia por canal (lib/analysis.js)},label={lst:analysis}]
function analyzeChannels(wifiHeatmap) {
  const scores = {};
  for (let ch = 1; ch <= 13; ch++) {
    let score = 0;
    for (let other = 1; other <= 13; other++) {
      const data = wifiHeatmap[other];
      if (!data || !data.count) continue;
      const dist = Math.abs(ch - other);
      if (dist > 4) continue;                        // sin solapamiento significativo
      const overlapWeight = 1 - dist / 5;            // decae con la distancia
      const rssiWeight = Math.max(0.2, Math.min(1, (data.maxRssi + 95) / 55));
      score += data.count * overlapWeight * rssiWeight;
    }
    scores[ch] = Math.round(score * 10) / 10;
  }
  // ... determina bestChannel, worstChannel y recommendedChannel (1/6/11)
  return { scores, bestChannel, worstChannel, recommendedChannel };
}
\end{lstlisting}

\subsection{Detección de gemelo malicioso}
La heurística agrupa por SSID y marca como sospechosos los que aparecen bajo más de un BSSID (Listado~\ref{lst:evil}).

\begin{lstlisting}[style=js,caption={Detección de evil twin (lib/analysis.js)},label={lst:evil}]
function detectEvilTwins(devices) {
  const bySsid = new Map();
  devices.forEach(d => {
    if (d.type !== 'wifi') return;
    const ssid = (d.ssid || '').trim();
    if (!ssid || ssid === 'Oculto') return;
    if (!bySsid.has(ssid)) bySsid.set(ssid, new Set());
    bySsid.get(ssid).add((d.mac || '').toLowerCase());
  });
  const flagged = new Set();
  bySsid.forEach(macs => { if (macs.size > 1) macs.forEach(m => flagged.add(m)); });
  return flagged;                                    // BSSIDs bajo el mismo SSID
}
\end{lstlisting}

\subsection{Identificación de fabricante (OUI)}
El módulo \code{lib/oui} mantiene un catálogo de 123 prefijos y detecta MAC aleatorizadas por el segundo \emph{nibble} del primer octeto (Listado~\ref{lst:oui}).

\begin{lstlisting}[style=js,caption={Búsqueda de fabricante y MAC aleatoria (lib/oui.js)},label={lst:oui}]
function lookupVendor(mac) {
  const clean = normalizeMac(mac);
  const prefix = clean.slice(0, 6);
  if (OUI_MAP[prefix]) return OUI_MAP[prefix];
  const secondNibble = clean[1];                     // MAC localmente administrada
  if (['2','6','A','E'].includes(secondNibble)) return 'MAC Aleatoria (privacidad)';
  return 'Desconocido';
}
\end{lstlisting}

\subsection{Autenticación por token}
El módulo \code{lib/auth} implementa sesiones por token con expiración y un \emph{middleware} \code{requireAuth} para proteger las rutas de control, además de \code{deviceAuthorized} para la ingesta (Listado~\ref{lst:auth}).

\begin{lstlisting}[style=js,caption={Sesiones y middleware de autorización (lib/auth.js)},label={lst:auth}]
function login(password) {
  if (password !== PASSWORD) return null;
  const token = crypto.randomBytes(24).toString('hex');
  tokens.set(token, Date.now() + TTL_MS);            // TTL 24 h
  return token;
}
function requireAuth(req, res, next) {
  const token = req.headers['x-eilor-token'] || req.query.token || (req.body && req.body.token);
  if (isValidToken(token)) return next();
  return res.status(401).json({ error: 'No autorizado. Inicia sesion.' });
}
function deviceAuthorized(t) { return !DEVICE_TOKEN || t === DEVICE_TOKEN; }
\end{lstlisting}

\subsection{Ingesta y difusión en tiempo real}
Al recibir un escaneo, \code{server.js} anota cada dispositivo con fabricante, MAC aleatoria y bandera de gemelo malicioso, recalcula el mapa de calor y difunde el resultado a los clientes (Listado~\ref{lst:ingest}).

\begin{lstlisting}[style=js,caption={Anotación y difusión (server.js)},label={lst:ingest}]
function annotateDevices(list) {
  return list.map(dev => ({
    ...dev,
    vendor: dev.vendor || lookupVendor(dev.mac),
    randomMac: dev.randomMac !== undefined ? dev.randomMac : isRandomMac(dev.mac),
    evilTwin: evilTwinMacs.has((dev.mac || '').toLowerCase())
  }));
}
// tras procesar el lote:
broadcast({ event: 'scan_update', heatmap: channelHeatmap,
            latestBatch: annotateDevices(Array.from(activeDevices.values())) });
\end{lstlisting}

\section{Implementación del frontend (dashboard)}
El cliente web establece la conexión WebSocket, solicita el estado inicial y despacha los eventos entrantes hacia los distintos renderizadores (mapa de calor, radar, tabla de dispositivos, registros y serie temporal). El Listado~\ref{lst:front} ilustra el manejo del evento de escaneo.

\begin{lstlisting}[style=js,caption={Manejo de eventos en el cliente (public/app.js)},label={lst:front}]
if (data.event === 'initial_state' || data.event === 'scan_update') {
  if (data.latestBatch) state.devices = data.latestBatch;
  if (data.heatmap)    { state.heatmap = data.heatmap; renderHeatmap(data.heatmap); }
  renderDevicesTable();
  renderRadarBlips();
  updateCounters();
}
\end{lstlisting}

\captura[0.98]{cap-espectro-radar}{Dashboard web en operación: mapa de saturación espectral Wi\nobreakdash-Fi (canales 1--13) y radar de proximidad BLE, alimentados en vivo por WebSocket.}
